To generate figures, some tricks may come in handy. For example,
\figref{plot} was generated in Matlab using
\begin{verbatim}
>> xlabel('x [mm]','interpreter','latex')
>> ylabel('y [mm]','interpreter','latex')
>> legend('Method 1','Method 2','location','SE','interpreter','latex')
\end{verbatim}
In order to get nice text in the figure.

Further, a figure can usually be cleaned up a bit by fixing the axis
and additionally scaling it to the final size:
\begin{verbatim}
>> set(gca,'xtick',[0 1],'ytick',1:3)
>> set(gcf,'Position',[100 100 240 180])
\end{verbatim}
This can be done similarly in other programming languages such as
Python and avoids scaling inside \LaTeX\ such that, e.g., the text
looks weirdly scaled.

\begin{figure}[htp]
  \centering
  \includegraphics{plot}
  \label{fig:plot}
  \caption{Comparision of Methods 1 and 2. \blabla{explanation goes
      here} \DO{Remind the reader of the meaning of the important
      variables. Consider to use different markers and not only colors
      for clarity (here: squares and circles, respectively. You should
      also refer to the figure in the main text.}  \DONT{Never insert
      a title in the figure, use the caption instead.}}
\end{figure}